\section*{अध्याय \ref{ch:b2:huckel} --- ह्युकेल सिद्धांत और संयुग्मित निकाय}

\begin{solution}{exo:b2:huckel:1}
ऐलिल स्तर ($n = 3$): $\alpha + 1.414\beta$, $\alpha$, $\alpha - 1.414\beta$। धनायन (2 इलेक्ट्रॉन):
$2\alpha + 2.828\beta$; मूलक (3): $3\alpha + 2.828\beta$; ऋणायन (4): $4\alpha + 2.828\beta$। अतिरिक्त
इलेक्ट्रॉन अनाबंधी स्तर में जाते हैं और केवल $\alpha$ भाग बदलते हैं।
\end{solution}

\begin{solution}{exo:b2:huckel:2}
$2(\alpha + 1.618\beta) + 2(\alpha + 0.618\beta) = 4\alpha + 4.472\beta$।
\end{solution}

\begin{solution}{exo:b2:huckel:3}\emergencystretch=3em
बेन्ज़ीन (6), साइक्लोपेंटाडाईनिल ऋणायन (6), ट्रोपिलियम (6): \omterm{def:b2:huckel:aromatic}{ऐरोमैटिक}। साइक्लोपेंटाडाईनिल
धनायन (4), साइक्लोब्यूटाडाईन (4), समतलीय साइक्लो-ऑक्टाटेट्राईन (8): समतलीय होने पर \omterm{def:b2:huckel:aromatic}{प्रति-ऐरोमैटिक}
(वास्तविक साइक्लो-ऑक्टाटेट्राईन टब-आकार का और अनैरोमैटिक है)।
\end{solution}

\begin{solution}{exo:b2:huckel:4}
वृत्त में एक अष्टभुज: स्तर $\alpha + 2\beta$, $\alpha + 1.414\beta$ (दो बार), $\alpha$ (दो बार),
$\alpha - 1.414\beta$ (दो बार), $\alpha - 2\beta$। आठ इलेक्ट्रॉन: दो निम्नतम में, चार
पहले युग्म में, एक-एक अनाबंधी युग्म के हर कक्षक में।
\end{solution}

\begin{solution}{exo:b2:huckel:5}
$4.472\beta - 4\beta = 0.472\beta$, अर्थात् $0.472 \times 75 = \qty{35}{kJ/mol}$, जबकि हाइड्रोजनन
आँकड़ों से लगभग \qty{10}{kJ/mol}: बेन्ज़ीन पर समंजित $\beta$ वाला ह्युकेल मॉडल
खुली शृंखलाओं के संयुग्मन को बढ़ा-चढ़ाकर आँकता है।
\end{solution}

\begin{solution}{exo:b2:huckel:6}
$p_{12} = p_{34} = 0.894$, $p_{23} = 0.447$: सिरे के आबंध C1--C2 और C3--C4
सबसे छोटे हैं, केंद्रीय आबंध लंबा परंतु एकल आबंध से छोटा।
\end{solution}

\begin{solution}{exo:b2:huckel:7}
ऋणायन, छह इलेक्ट्रॉन: $2 \times 2 + 4 \times 0.618$, $E_\pi = 6\alpha + 6.472\beta$, संवृत कोश:
\omterm{def:b2:huckel:aromatic}{ऐरोमैटिक}। धनायन, चार: $2 \times 2 + 2 \times 0.618$, $4\alpha + 5.236\beta$, अपभ्रष्ट युग्म के
हर कक्षक में एक इलेक्ट्रॉन: \omterm{def:b2:huckel:aromatic}{प्रति-ऐरोमैटिक}, अत्यंत अस्थायी।
\end{solution}

\begin{solution}{exo:b2:huckel:8}
सप्तसदस्यी वलय: $\alpha + 2\beta$, $\alpha + 1.247\beta$ (दो बार), फिर प्रतिआबंधी स्तर। छह
इलेक्ट्रॉन, जैसे धनायन \ce{C7H7+} में, आबंधी स्तरों को ठीक-ठीक भरते हैं: एक \omterm{def:b2:huckel:aromatic}{ऐरोमैटिक}
धनायन। यौगिक आसानी से ट्रोपिलियम और ब्रोमाइड में आयनित होता है।
\end{solution}

\begin{solution}{exo:b2:huckel:9}
$1.802 + 1.247 + 0.445 - 0.445 - 1.247 - 1.802 = 0$: ऊर्जाओं का योग $6\alpha$ है।
\end{solution}

\begin{solution}{exo:b2:huckel:10}
$n = 6$: $m_k = 2\cos(k\pi/7) = 1.802$, $1.247$, $0.445$, $-0.445$, $-1.247$, $-1.802$। छह
इलेक्ट्रॉन: $E_\pi = 6\alpha + 2(1.802 + 1.247 + 0.445)\beta = 6\alpha + 6.988\beta$;
विस्थानीकरण $0.988\beta$।
\end{solution}

\begin{solution}{exo:b2:huckel:11}
वर्ग में उच्चतम ऊर्जा वाले दो इलेक्ट्रॉन एक अपभ्रष्ट अनाबंधी युग्म के हर कक्षक में
एक-एक बैठते हैं। आमने-सामने के दो आबंधों को खींचने से अपभ्रष्टता हट जाती है:
एक कक्षक नीचे जाता है, दूसरा ऊपर, और दोनों इलेक्ट्रॉन निचले में युग्मित हो जाते हैं, जिससे
ऊर्जा घटती है। अणु दो स्थानीकृत द्वि-आबंधों वाले आयत के रूप में स्थिर होता है;
वह \omterm{def:b2:huckel:aromatic}{प्रति-ऐरोमैटिक} और अत्यंत अभिक्रियाशील बना रहता है।
\end{solution}

\begin{solution}{exo:b2:huckel:12}
आव्यूह ($\beta$ की इकाइयों में, $\alpha$ को मूल मानकर) $\begin{pmatrix}0 & 1\\ 1 & 1\end{pmatrix}$:
$m = 1.618$ और $-0.618$, $E = \alpha + 1.618\beta$ (आबंधी) और $\alpha - 0.618\beta$।
\omterm{def:b2:diatomic-mos:bonding}{आबंधी कक्षक} के लिए $c_X = 1.618c_C$, अतः $c_C^2 = 0.276$, $c_X^2 = 0.724$; दो इलेक्ट्रॉनों के साथ
$q_C = 0.55$, $q_X = 1.45$: $\pi$ घनत्व X की ओर खिंचता है।
\end{solution}

\begin{solution}{pb:b2:huckel:1}
\textbf{1.} $-123.1 - (-4.3) = \qty{-118.8}{kJ/mol}$।
\textbf{2.} $-123.1 - 104.6 = \qty{-227.7}{kJ/mol}$।
\textbf{3.} $-123.1 - 82.9 = \qty{-206.0}{kJ/mol}$।
\textbf{4.} $3 \times (-118.8) = \qty{-356.4}{kJ/mol}$।
\textbf{5.} $356.4 - 206.0 = \qty{150}{kJ/mol}$।
\textbf{6.} डाईन दो पृथक द्वि-आबंधों से $2 \times 118.8 - 227.7 = \qty{10}{kJ/mol}$ कम
मुक्त करता है: खुली शृंखला का संयुग्मन छोटा लाभ देता है, बेन्ज़ीन का बंद
वलय पंद्रह गुना अधिक।
\textbf{7.} $6 \times 6$ आव्यूह, जिसमें हर परमाणु और उसके दोनों पड़ोसियों के बीच 1 (1--2, 2--3,
\dots, 6--1) और अन्यत्र 0 है।
\textbf{8.} $m = 2\cos(2\pi k/6)$: $2$, $1$, $1$, $-1$, $-1$, $-2$: $E = \alpha + 2\beta$, $\alpha + \beta$ (दो बार),
$\alpha - \beta$ (दो बार), $\alpha - 2\beta$।
\textbf{9.} दो इलेक्ट्रॉन $\alpha + 2\beta$ में, चार युग्म $\alpha + \beta$ में।
\textbf{10.} $E_\pi = 6\alpha + 8\beta$।
\textbf{11.} $3(2\alpha + 2\beta) = 6\alpha + 6\beta$।
\textbf{12.} $2\beta$, अर्थात् $2|\beta|$ का स्थायीकरण।
\textbf{13.} $2 + 1 + 1 - 1 - 1 - 2 = 0$: छहों स्तरों का योग $6\alpha$ है।
\textbf{14.} $2|\beta| = \qty{150}{kJ/mol}$: $|\beta| = \qty{75}{kJ/mol}$।
\textbf{15.} $0.472 \times 75 = \qty{35}{kJ/mol}$, जबकि \qty{10}{kJ/mol}: बेन्ज़ीन पर समंजित
मॉडल खुली शृंखला को बढ़ा-चढ़ाकर आँकता है।
\textbf{16.} भरे कक्षक $\psi_1 = (1, 1, 1, 1, 1, 1)/\sqrt6$,
$\psi_2 = (2, 1, -1, -2, -1, 1)/\sqrt{12}$ और $\psi_3 = (0, 1, 1, 0, -1, -1)/2$ लिए जा सकते हैं। आबंध
1--2 के लिए: $p_{12} = 2(\frac16 + \frac{2}{12} + 0) = \frac23$; सममिति से हर आबंध का $p = 2/3$ है।
\textbf{17.} ब्यूटाडाईन के सिरे के (0.894) और केंद्रीय (0.447) आबंधों के बीच।
\textbf{18.} छहों आबंध सममिति से तुल्य हैं; उनकी लंबाई, \qty{139.7}{pm},
एथीन के द्वि-आबंध (\qty{133.9}{pm}) और एथेन के एकल आबंध
(\qty{153.6}{pm}) के बीच है।
\textbf{19.} $0.988|\beta| = \qty{74}{kJ/mol}$, बेन्ज़ीन का आधा: वलय को बंद करना
\omterm{def:b2:huckel:delocalisation}{विस्थानीकरण ऊर्जा} को दुगुना कर देता है।
\textbf{20.} $\alpha + 2\beta$, $\alpha + 1.414\beta$ (दो बार), $\alpha$ (दो बार), $\alpha - 1.414\beta$ (दो बार),
$\alpha - 2\beta$।
\textbf{21.} अंतिम दो इलेक्ट्रॉन अनाबंधी युग्म के हर कक्षक में एक-एक जाते हैं: खुला
कोश, \omterm{def:b2:huckel:aromatic}{प्रति-ऐरोमैटिक}।
\textbf{22.} यह मुड़कर टब बन जाता है: पड़ोसी द्वि-आबंधों के p कक्षक अब
समांतर नहीं रहते, आबंध एकांतर होते हैं, और अणु पॉलिईन की तरह व्यवहार करता है।
\textbf{23.} चार इलेक्ट्रॉन: वही आधा भरा युग्म, और भी बुरा क्योंकि अणु समतलता से
बच नहीं सकता; वह आयत में विकृत हो जाता है और अत्यंत अभिक्रियाशील है।
\textbf{24.} समतलीय एकचक्रीय संयुग्मित निकाय $4N + 2$ $\pi$ इलेक्ट्रॉनों के साथ \omterm{def:b2:huckel:aromatic}{ऐरोमैटिक} है,
$4N$ के साथ \omterm{def:b2:huckel:aromatic}{प्रति-ऐरोमैटिक}।
\textbf{25.} $\boldsymbol{|\beta| \approx \qty{75}{kJ/mol}}$।
\end{solution}
